\documentclass[12pt,a4paper]{article}

\usepackage[utf8]{inputenc}
\usepackage[T1]{fontenc}
\usepackage[french]{babel}
\usepackage{geometry}
\usepackage{graphicx}
\usepackage{booktabs}
\usepackage{array}
\usepackage{xcolor}
\usepackage{listings}
\usepackage{hyperref}
\usepackage{fancyhdr}
\usepackage{titlesec}
\usepackage{enumitem}
\usepackage{tabularx}
\usepackage{multirow}
\usepackage{float}

\geometry{top=2.5cm, bottom=2.5cm, left=2.5cm, right=2.5cm}

% Couleurs
\definecolor{petblue}{RGB}{37,99,235}
\definecolor{codegreen}{RGB}{16,185,129}
\definecolor{codegray}{RGB}{107,114,128}
\definecolor{lightgray}{RGB}{243,244,246}
\definecolor{darkbg}{RGB}{17,24,39}

% Listings
\lstdefinestyle{codestyle}{
  backgroundcolor=\color{lightgray},
  commentstyle=\color{codegray},
  keywordstyle=\color{petblue}\bfseries,
  stringstyle=\color{codegreen},
  basicstyle=\ttfamily\small,
  breakatwhitespace=false,
  breaklines=true,
  frame=single,
  rulecolor=\color{codegray},
  keepspaces=true,
  numbers=left,
  numberstyle=\tiny\color{codegray},
  showspaces=false,
  showstringspaces=false,
  tabsize=2
}
\lstset{style=codestyle}

% En-tête / pied de page
\pagestyle{fancy}
\fancyhf{}
\rhead{\textcolor{petblue}{\textbf{PET Platform}}}
\lhead{Rapport de développement}
\rfoot{\thepage}
\lfoot{\textit{Abdelmoumen Elimrany}}

% Titres
\titleformat{\section}{\Large\bfseries\color{petblue}}{}{0em}{}[\titlerule]
\titleformat{\subsection}{\large\bfseries\color{darkbg}}{}{0em}{}

% Liens
\hypersetup{
  colorlinks=true,
  linkcolor=petblue,
  urlcolor=petblue,
  pdftitle={Rapport PET Platform},
  pdfauthor={Abdelmoumen Elimrany}
}

% ─────────────────────────────────────────────────────────────────────────────

\begin{document}

% Page de titre
\begin{titlepage}
  \centering
  \vspace*{2cm}

  {\Huge\bfseries\color{petblue} PET Platform}\\[0.5cm]
  {\Large Rapport de développement}\\[0.3cm]
  {\large Privacy-Enhancing Technologies — Plateforme de soumission sécurisée}

  \vspace{1.5cm}
  \hrule height 1pt
  \vspace{0.5cm}

  {\large\bfseries Implémentation des issues \#6 et \#7}\\[0.3cm]
  {\normalsize Support R / Collecte d'outputs / Navigateur sécurisé Brave}

  \vspace{0.5cm}
  \hrule height 1pt

  \vfill

  \begin{tabular}{rl}
    \textbf{Étudiant :} & Abdelmoumen Elimrany \\
    \textbf{GitHub :}   & \url{https://github.com/ElimranyAbdelmoumen/PET_Projet} \\
    \textbf{Date :}     & 22 avril 2026 \\
    \textbf{Branche :}  & \texttt{main} \\
  \end{tabular}

  \vspace{2cm}
\end{titlepage}

\tableofcontents
\newpage

% ─────────────────────────────────────────────────────────────────────────────
\section{Présentation du projet}

La \textbf{PET Platform} (Privacy-Enhancing Technologies) est une application web permettant à des chercheurs de soumettre des scripts Python ou R pour exécution dans un environnement isolé. Les résultats sont vérifiés par un administrateur avant d'être accessibles à l'utilisateur.

\subsection{Architecture générale}

\begin{center}
\begin{tabular}{|l|l|l|}
\hline
\textbf{Composant} & \textbf{Technologie} & \textbf{Rôle} \\
\hline
Backend & Flask (Python) & Interface web, authentification, API \\
Base de données & PostgreSQL 16 & Stockage des soumissions et outputs \\
Worker & Python + Docker SDK & Exécution asynchrone des scripts \\
Sandbox & Docker (isolé) & Conteneur d'exécution sécurisé \\
Navigateur & Brave + Extension & Accès sécurisé à la plateforme \\
\hline
\end{tabular}
\end{center}

\vspace{0.5cm}

Le worker scrute la base de données toutes les 3 secondes, récupère les soumissions avec le statut \texttt{APPROVED}, les exécute dans un conteneur Docker éphémère, puis enregistre les résultats.

% ─────────────────────────────────────────────────────────────────────────────
\section{Issue \#6 — Support R, sélecteur de fichiers, collecte des outputs}

\subsection{Objectif}

L'issue demandait trois améliorations majeures de la plateforme :
\begin{enumerate}
  \item Support de l'exécution de scripts \textbf{R} en plus de Python
  \item \textbf{Sélecteur de fichier microdata} pour associer un jeu de données à une soumission
  \item \textbf{Collecte et validation des fichiers de sortie} produits par les scripts
\end{enumerate}

% ──────────────────────────────────────────────
\subsection{Partie 1 — Support R}

\subsubsection*{Modification du runner Docker}

Le fichier \texttt{runner/Dockerfile} a été mis à jour pour inclure \texttt{r-base} et \texttt{r-base-dev} (nécessaire pour compiler \texttt{rpy2}) :

\begin{lstlisting}[language=bash, caption={runner/Dockerfile}]
RUN apt-get update && apt-get install -y --no-install-recommends \
    r-base \
    r-base-dev \
    && rm -rf /var/lib/apt/lists/*

RUN pip install --no-cache-dir \
    pandas numpy matplotlib scikit-learn scipy \
    seaborn requests openpyxl xlrd rpy2

RUN mkdir -p /work/outputs /work/data && chmod 777 /work/outputs /work/data
ENTRYPOINT ["python", "/usr/local/bin/run_script.py"]
\end{lstlisting}

\textbf{Note :} L'utilisateur \texttt{runner} (UID 10001) a été supprimé car il causait une erreur de permissions lors de l'écriture dans \texttt{/work/outputs/}.

\subsubsection*{Détection de l'extension dans le worker}

Le worker (\texttt{runner\_worker/worker.py}) détecte l'extension du script pour choisir l'interpréteur :

\begin{lstlisting}[language=Python, caption={Détection Python/R dans worker.py}]
ext = os.path.splitext(file_path)[1].lower() or ".py"
container_script = f"/work/script{ext}"
# ext = ".py"  => python run_script.py /work/script.py
# ext = ".R"   => Rscript --vanilla /work/script.R
\end{lstlisting}

Le script wrapper \texttt{run\_script.py} route l'exécution selon l'extension :

\begin{lstlisting}[language=Python, caption={run\_script.py — détection Python/R}]
if script_path.endswith(".R"):
    result = subprocess.run(
        ["Rscript", "--vanilla", script_path],
        capture_output=True, text=True, timeout=55
    )
else:
    # exec() Python avec capture des variables
    ...
\end{lstlisting}

% ──────────────────────────────────────────────
\subsection{Partie 2 — Sélecteur de fichier microdata}

La table \texttt{microdata\_files} stocke les jeux de données disponibles (CSV, Excel). L'utilisateur peut sélectionner un fichier microdata lors de la soumission de son script.

\subsubsection*{Schéma de base de données}

\begin{lstlisting}[language=SQL, caption={Tables microdata\_files et submissions (extrait)}]
CREATE TABLE microdata_files (
    id        SERIAL PRIMARY KEY,
    guid      UUID DEFAULT gen_random_uuid() UNIQUE,
    name      VARCHAR(255) NOT NULL,
    file_path TEXT NOT NULL,
    created_at TIMESTAMPTZ DEFAULT NOW()
);

ALTER TABLE submissions
    ADD COLUMN microdata_guid UUID REFERENCES microdata_files(guid),
    ADD COLUMN outputs_status VARCHAR(20) DEFAULT 'NONE';
\end{lstlisting}

Le fichier microdata est monté en lecture seule dans le conteneur sous \texttt{/work/data/<nom\_fichier>}.

% ──────────────────────────────────────────────
\subsection{Partie 3 — Collecte et validation des outputs}

\subsubsection*{Flux de validation}

\begin{center}
\begin{tabular}{|c|l|}
\hline
\textbf{Étape} & \textbf{Description} \\
\hline
1 & Script écrit ses fichiers dans \texttt{/work/outputs/} \\
2 & Worker scanne le répertoire local après exécution \\
3 & Fichiers insérés dans \texttt{submission\_outputs} (statut \texttt{PENDING\_VALIDATION}) \\
4 & Admin visualise et approuve ou rejette les outputs \\
5 & Utilisateur télécharge ses fichiers si \texttt{APPROVED} \\
\hline
\end{tabular}
\end{center}

\subsubsection*{Correction du bug chemin local/hôte}

Le worker utilise deux chemins distincts :

\begin{lstlisting}[language=Python, caption={Séparation chemin hôte vs chemin local}]
# Pour docker run -v (chemin vu par Docker Desktop)
host_outputs_dir = os.path.join(HOST_STORAGE_PATH, "outputs", str(submission_id))

# Pour os.listdir (chemin dans le container worker)
local_outputs_dir = os.path.join(STORAGE_PATH, "outputs", str(submission_id))

# HOST_STORAGE_PATH = /host_mnt/c/Users/User/Desktop/PET_Projet/storage
# STORAGE_PATH      = /storage  (monté par docker-compose)
\end{lstlisting}

\subsubsection*{Table submission\_outputs}

\begin{lstlisting}[language=SQL]
CREATE TABLE submission_outputs (
    id            SERIAL PRIMARY KEY,
    submission_id INTEGER REFERENCES submissions(id),
    filename      VARCHAR(255) NOT NULL,
    file_path     TEXT NOT NULL,
    created_at    TIMESTAMPTZ DEFAULT NOW()
);
\end{lstlisting}

% ─────────────────────────────────────────────────────────────────────────────
\section{Issue \#7 — Navigateur sécurisé Brave}

\subsection{Objectif}

Restreindre l'accès à la plateforme au seul navigateur \textbf{Brave} avec l'extension \textbf{Brave Security Shield} installée, et activer les protections suivantes :

\begin{itemize}
  \item Blocage du copier/coller
  \item Blocage des captures d'écran (PrintScreen)
  \item Blocage du partage d'écran (\texttt{getDisplayMedia})
  \item Blocage de l'inspection (DevTools / F12)
  \item Filigrane (watermark) avec l'identité de l'utilisateur connecté
  \item Journal des violations
\end{itemize}

\subsection{Architecture de l'extension}

L'extension est un \textbf{Manifest V3} pour Brave/Chromium composée de :

\begin{center}
\begin{tabular}{|l|l|}
\hline
\textbf{Fichier} & \textbf{Rôle} \\
\hline
\texttt{manifest.json}      & Déclaration de l'extension (MV3, permissions) \\
\texttt{content\_script.js} & Protections injectées dans chaque page \\
\texttt{background.js}      & Service worker — stockage violations, identité \\
\texttt{popup.html/js}      & Interface admin (toggles, watermark, journal) \\
\texttt{screen\_block.css}  & Overlay CSS de blocage visuel \\
\texttt{auth\_bridge.js}    & Pont d'authentification avec Flask \\
\hline
\end{tabular}
\end{center}

\subsection{Intégration Flask}

\subsubsection*{Mécanisme de détection}

\texttt{content\_script.js} pose un cookie \texttt{bss\_active=1} à chaque visite sur la plateforme. Flask vérifie ce cookie sur toutes les routes via le décorateur \texttt{browser\_required} :

\begin{lstlisting}[language=Python, caption={authz.py — décorateur browser\_required}]
def browser_required(fn):
    @wraps(fn)
    def wrapper(*args, **kwargs):
        if not request.cookies.get("bss_active"):
            return render_template("browser_required.html"), 403
        return fn(*args, **kwargs)
    return wrapper
\end{lstlisting}

Sans extension $\rightarrow$ HTTP 403 avec instructions d'installation.\\
Avec extension $\rightarrow$ accès normal.

\subsubsection*{Initialisation du watermark}

Après connexion, chaque page protégée transmet l'identité Flask à l'extension via \texttt{auth\_bridge.js} :

\begin{lstlisting}[language=HTML, caption={Injection de l'identité dans les templates}]
<script>
  var BSS_USER_ID = {{ session.get('user_id') }};
  var BSS_USERNAME = {{ session.get('username') | tojson }};
</script>
<script type="module" src="/static/js/bss_init.js"></script>
\end{lstlisting}

\begin{lstlisting}[language=JavaScript, caption={bss\_init.js}]
import { initBSSAuth } from "./auth_bridge.js";
if (BSS_USER_ID) {
  initBSSAuth({ uid: String(BSS_USER_ID), displayName: BSS_USERNAME });
}
\end{lstlisting}

Le watermark affiche alors le nom de l'utilisateur connecté en filigrane diagonal sur toutes les pages.

\subsection{Installation de l'extension}

\begin{enumerate}
  \item Installer \textbf{Brave Browser} (\url{https://brave.com})
  \item Ouvrir \texttt{brave://extensions} et activer le \textbf{mode développeur}
  \item Cliquer sur \textit{Charger l'extension non empaquetée}
  \item Sélectionner le dossier \texttt{breavescripts/} du projet
  \item Copier l'\textbf{Extension ID} et le remplacer dans \texttt{auth\_bridge.js} ligne 37
\end{enumerate}

% ─────────────────────────────────────────────────────────────────────────────
\section{Benchmark du sandbox}

\subsection{Conditions d'exécution}

\begin{itemize}
  \item \textbf{Image :} \texttt{portwatch-python-runner:1.0}
  \item \textbf{Limites :} 512 Mo RAM, 1 CPU, réseau désactivé (\texttt{--network none})
  \item \textbf{Timeout :} 60 secondes
  \item \textbf{Méthode :} 5 exécutions par scénario, mesure du temps total docker run
  \item \textbf{Machine :} Windows 11, Docker Desktop (WSL2 backend)
\end{itemize}

\subsection{Résultats}

\begin{table}[H]
\centering
\caption{Temps d'exécution du sandbox (ms) — 5 runs par scénario}
\begin{tabularx}{\textwidth}{|X|c|c|c|c|c|c|}
\hline
\textbf{Scénario} & \textbf{R1} & \textbf{R2} & \textbf{R3} & \textbf{R4} & \textbf{R5} & \textbf{Moy.} \\
\hline
Python — hello world           & 1611 & 2070 & 1417 & 1689 & 1379 & \textbf{1633} \\
Python — pandas describe       & 1347 & 1502 & 1351 & 1271 & 1361 & \textbf{1366} \\
Python — sklearn LinearReg.    & 1395 & 1514 & 1253 & 1484 & 1427 & \textbf{1414} \\
Python — lecture microdata CSV & 1293 & 1593 & 1395 & 1545 & 1737 & \textbf{1512} \\
Python — écriture output CSV   & 1570 & 1601 & 1501 & 1344 & 1398 & \textbf{1482} \\
\hline
R — hello world                & 1613 & 1613 & 1375 & 1423 & 1344 & \textbf{1473} \\
R — stats rnorm(1000)          & 1743 & 1483 & 1339 & 1600 & 1484 & \textbf{1529} \\
R — lecture microdata CSV      & 1424 & 1759 & 1631 & 1464 & 1332 & \textbf{1522} \\
\hline
\end{tabularx}
\end{table}

\subsection{Analyse}

\textbf{Temps de démarrage du conteneur.}
Le temps dominant (\textasciitilde 1,3–1,7 s) correspond au démarrage du conteneur Docker éphémère, pas au calcul. Le code lui-même s'exécute en quelques millisecondes.

\textbf{Python vs R.}
Les deux langages ont des performances quasi identiques dans ce sandbox (différence $<$ 200 ms). La variabilité observée est liée à la charge système, pas à l'interpréteur.

\textbf{Overhead microdata et outputs.}
L'ajout d'un fichier microdata ou d'une écriture de fichier output ajoute moins de 150 ms — overhead négligeable.

\textbf{Stabilité.}
Aucun timeout ni erreur sur 40 exécutions. L'écart-type estimé est de ±150 ms, cohérent avec la variabilité du démarrage Docker.

\subsection{Synthèse des contraintes}

\begin{center}
\begin{tabular}{|l|l|c|}
\hline
\textbf{Contrainte} & \textbf{Valeur} & \textbf{Respectée} \\
\hline
Isolation réseau   & \texttt{--network none}    & \checkmark \\
Mémoire max        & 512 Mo                     & \checkmark \\
CPU max            & 1 core                     & \checkmark \\
Timeout            & 60 s (moy. \textasciitilde 1,5 s) & \checkmark \\
Support Python     & pandas, numpy, sklearn...  & \checkmark \\
Support R          & stats, utils, base         & \checkmark \\
Lecture microdata  & \texttt{/work/data/}       & \checkmark \\
Écriture outputs   & \texttt{/work/outputs/}    & \checkmark \\
\hline
\end{tabular}
\end{center}

% ─────────────────────────────────────────────────────────────────────────────
\section{Conclusion}

Les deux issues demandées par l'encadrant ont été implémentées et mergées sur la branche \texttt{main} :

\begin{itemize}
  \item \textbf{Issue \#6} — Support complet R + sélecteur microdata + collecte/validation des outputs avec workflow admin
  \item \textbf{Issue \#7} — Intégration de l'extension Brave Security Shield : accès interdit sans l'extension (HTTP 403), watermark utilisateur, blocage copier/coller, screenshots et DevTools
\end{itemize}

Le sandbox Docker isole correctement chaque soumission. Le coût fixe de démarrage (\textasciitilde 1,5 s) est largement acceptable pour une plateforme de traitement de code scientifique.

\vspace{1cm}
\begin{flushright}
\textit{Abdelmoumen Elimrany} \\
\textit{22 avril 2026}
\end{flushright}

\end{document}
